%(BEGIN_QUESTION)
% Copyright 2006, Tony R. Kuphaldt, released under the Creative Commons Attribution License (v 1.0)
% This means you may do almost anything with this work of mine, so long as you give me proper credit

Anta at du skal konvertere lengden 35 fot til tommer. Utregningen er enkel, men du ønsker å øve på bruk av «enhetsbrøker»:

$${{35 \hbox{ ft}} \over {1}} \times {{12 \hbox{ in}} \over {1 \hbox{ ft}}}$$

Enheten «fot» (ft) kanselleres, og du får svaret 420 tommer. Så langt, så bra.

\vskip 10pt

Anta nå at du må konvertere arealet 70 kvadratfot (ft$^{2}$) til kvadrattommer (in$^{2}$). Det må finnes en måte å bruke «enhetsbrøk»-konseptet på her også, men det fungerer ikke hvis vi setter det opp identisk med forrige oppgave:

$${{70 \hbox{ ft}^2} \over {1}} \times {{12 \hbox{ in}} \over {1 \hbox{ ft}}} = \hbox{ Feil svar!}$$

Grunnen til at dette ikke fungerer som endelig løsning på arealproblemet er at kvadratfot ikke kanselleres med lineære fot. Størrelsene er fortsatt én-dimensjonale: fot, ikke kvadratfot; tommer, ikke kvadrattommer.

Men hva om vi gjør dette?

$${{70 \hbox{ ft}^2} \over {1}} \times {{12^2 \hbox{ in}^2} \over {1^2 \hbox{ ft}^2}}$$

... som er det samme som dette ...

$${{70 \hbox{ ft}^2} \over {1}} \times {{12 \hbox{ in}} \over {1 \hbox{ ft}}} \times {{12 \hbox{ in}} \over {1 \hbox{ ft}}}$$

Gir dette riktig svar? Hvorfor eller hvorfor ikke? Hvis ikke, forklar hvordan vi {\it kan} få riktig svar ved bruk av enhetsbrøkmetoden.

Til slutt: når du har bestemt en korrekt metode for arealenheter, generaliser regelen slik at den også dekker omregning av volumenheter ({\it kubikk}fot og tommer -- ft$^{3}$ og in$^{3}$).

\underbar{file i00107}
%(END_QUESTION)





%(BEGIN_ANSWER)

70 ft$^{2}$ = 10,080 in$^{2}$

\vskip 10pt

Det er mulig å ta en hvilken som helst en-dimensjonal enhetsbrøk og gjøre den om til en tilsvarende to- eller tre-dimensjonal enhetsbrøk ved å opphøye hele brøken i en potens. Siden en enhetsbrøk alltid har fysisk verdi 1, vil kvadrering eller kubering av den ikke endre verdien, like lite som å opphøye $1 \over 1$ i andre eller tredje potens.

Noen eksempler:

$$\left({{12 \hbox{ in}} \over {1 \hbox{ ft}}}\right)^2 = {{144 \hbox{ in}^2} \over {1 \hbox{ ft}^2}}$$

\vskip 10pt

$$\left({{12 \hbox{ in}} \over {1 \hbox{ ft}}}\right)^3 = {{1728 \hbox{ in}^3} \over {1 \hbox{ ft}^3}}$$

\vskip 10pt

$$\left({{1 \hbox{ in}} \over {2.54 \hbox{ cm}}}\right)^2 = {{1 \hbox{ in}^2} \over {6.4516 \hbox{ cm}^2}}$$

\vskip 10pt

$$\left({{1 \hbox{ in}} \over {2.54 \hbox{ cm}}}\right)^3 = {{1 \hbox{ in}^3} \over {16.387 \hbox{ cm}^3}}$$

Dette er nyttig, fordi det reduserer antall omregningsfaktorer du må huske. I stedet for å memorere 12 tommer = 1 fot {\it og} 144 kvadrattommer = 1 kvadratfot {\it og} 1728 kubikktommer = 1 kubikkfot, holder det å huske 12 tommer = 1 fot og så kvadrere eller kubere når det trengs.

%(END_ANSWER)





%(BEGIN_NOTES)


%INDEX% Physics, units and conversions: unit conversions using unity fractions

%(END_NOTES)
